\documentclass[11pt]{ctexart}
\usepackage[a4paper,margin=1in]{geometry}
\usepackage{graphicx}
\usepackage{xcolor}
\usepackage{hyperref}
\definecolor{refgray}{gray}{0.45}
\title{\textbf{Market Views｜全球投行叙事汇编}\\[0.4em]{\large 存储周期结构性分化与汽车电动化区域裂痕成为今日跨主题主线}}
\author{KC桌面 自动生成}
\date{260713}
\begin{document}
\maketitle
\tableofcontents
\newpage
\section{一页摘要}
\begin{itemize}
  \item 存储芯片行业共识AI需求强劲驱动DRAM/NAND出货，但分歧在于周期是否见顶：部分投行认为超级周期延续，逻辑从全面涨价转向结构性分化。
  \item 中国半导体资本开支上行周期获一致确认，但长鑫存储IPO对低端DRAM市场的竞争影响存在分歧。
  \item 欧洲电动化渗透率在平价车型推动下显著提升，EU-5六月同比增8.6个百分点至34.5\%，但中国新能源车高渗透率掩盖了乘用车总量同比降23\%的调整。
  \item 豪华车经销商新车亏损收窄仅为数学效应，售后业务因门店关闭恶化，行业整合尚未完成，分歧在于触底时点。
  \item 日本车企重新将人形机器人视为增长引擎，三菱量产计划（2027年月产1000台）被视作商业化信号，但安川电机利润率仅1.6\%，盈利转化存疑。
  \item 中国现制茶饮头部品牌同店销售环比改善，韧性超预期；正餐板块（海底捞、太二）6月同店增速放缓，暑期消费前景需重新校准。
  \item 中国电子纱/电子布供需紧张、价格高位，与粗纱受建筑淡季拖累形成分化；印度市场外资持仓降至极低水平，资金回流空间充足。
  \item 油轮运价受地缘事件扰动但核心航线运价同比大幅上涨，分歧在于三季度有效利用率恢复后运价可持续性。
\end{itemize}
\section{全球投行叙事汇编}
今日纳入 14 篇投行/券商研究，按 4 个市场、资产与产业主题系统整理。
\newpage
\subsection{存储与AI：周期延续与结构性分化}
\textbf{存储芯片行业正经历从全面涨价向结构性分化的转变，AI服务器需求成为核心驱动力，而中国半导体资本开支上行周期则为本土代工企业带来长期增长机遇。}
\subsubsection*{主流共识}
\begin{itemize}
  \item AI服务器和通用服务器需求强劲，持续拉动DRAM和NAND出货，是当前存储周期的核心驱动力。
  \item 存储芯片供给端约束具有黏性，长期协议模式正在缓解行业周期性波动。
  \item 中国半导体先进制程资本开支进入上行周期，本土代工企业扩产计划明确。
\end{itemize}
\subsubsection*{分歧与条件差异}
\begin{itemize}
  \item 市场担忧存储周期见顶，但部分机构认为当前是超级周期的延续，驱动逻辑从全面涨价转向结构性分化。
  \item 对长鑫存储IPO的影响存在分歧：部分观点认为其扩产是情绪扰动，实际冲击有限；另一些观点则关注其对低端DRAM市场的潜在竞争。
  \item 腾讯AI投入对短期利润的影响节奏存在分歧：JEF预计二季度非IFRS净利润增速仅4\%，远低于市场共识的9.7\%。
\end{itemize}
\subsubsection*{机构观点}
\paragraph{BofA} 存储芯片周期并非尾声，而是超级周期的延续，驱动逻辑从全面涨价转向结构性分化。自建景气指标5月录得183，远高于2017年和2024年周期高点（120-130）。DRAM和NAND全球出货金额同比增幅均超300\%，韩国半导体出口同比增199\%，DRAM现货价格同比涨幅达700\%。渠道调研显示美国大型科技公司需求依然强劲，长鑫存储实际冲击有限，三季度DRAM合约价已有多个新合同显示20\%以上涨幅，因此将三季度DRAM均价环比涨幅预测从17\%上调至21\%。
{\small\color{refgray}数据：BofA自建存储景气指标5月录得183，远高于历史周期高点120-130\par}
{\small\color{refgray}数据：DRAM和NAND全球出货金额同比增幅均超300\%\par}
{\small\color{refgray}数据：韩国半导体出口同比增199\%\par}
{\small\color{refgray}数据：DRAM现货价格同比涨幅达700\%\par}
{\small\color{refgray}边际变化：市场担忧存储周期见顶，但BofA基于自建指标和渠道调研，认为当前是超级周期的延续，并上调三季度ASP预测。\par}
\paragraph{GS} 南亚科技2Q26财报验证了DRAM行业的强劲涨价趋势，ASP连续两季环比涨幅超60\%（1Q26超70\%，2Q26超60\%），营业利润率从61\%跳升至74\%。AI和通用服务器正在驱动存储器需求增长，同时挤压非服务器领域供给，导致供需格局持续偏紧。GS据此推算三星电子常规DRAM在2Q26的ASP环比上涨约47\%，DRAM营业利润率达80\%，接近上一轮周期顶部水平。长期协议模式正在缓解存储行业的周期性波动。
{\small\color{refgray}数据：南亚科技2Q26 DRAM ASP环比涨超60\%，1Q26超70\%\par}
{\small\color{refgray}数据：南亚科技营业利润率从61\%跳升至74\%\par}
{\small\color{refgray}数据：GS推算三星电子2Q26常规DRAM ASP环比涨47\%，DRAM OPM达80\%\par}
{\small\color{refgray}数据：南亚科技预计供给紧张将持续数个季度\par}
{\small\color{refgray}边际变化：GS通过南亚科技财报读入，确认存储涨价周期持续，且涨价驱动力从消费电子转向AI服务器，供给紧张预期延长。\par}
\paragraph{GS} 晶合集成受益于中国半导体先进制程资本开支上行周期，40nm和28nm扩产计划明确。公司已宣布总预算355亿元的建厂方案，预计2026年四季度试产，2028年二季度满产，届时新增月产能5.5万片12英寸晶圆。GS将估值方法从近端P/E切换为2030年贴现P/E，反映长期可见度提升。上调2026-2027年收入预测1\%和8\%，但2027年毛利率因新产能折旧下调0.9个百分点至29.4\%。公司宣布2026年6月起执行10\%价格调整以对冲原材料成本。
{\small\color{refgray}数据：晶合集成40nm和28nm扩产总预算355亿元\par}
{\small\color{refgray}数据：新厂预计2026年四季度试产，2028年二季度满产，月产能5.5万片12英寸晶圆\par}
{\small\color{refgray}数据：GS上调2026-2027年收入预测1\%和8\%\par}
{\small\color{refgray}数据：2027年毛利率预测下调0.9个百分点至29.4\%\par}
{\small\color{refgray}边际变化：GS将估值方法从近端P/E切换为2030年贴现P/E，反映晶合集成从‘能否量产’进入‘规模化盈利’阶段，但新产能折旧压力上升。\par}
\paragraph{JEF} 腾讯AI产品线全面铺开，但投入加大导致短期利润承压。预计二季度非IFRS净利润增速仅4\%，远低于市场共识的9.7\%。微信AI‘小微’6月启动小范围测试，正式版预计4Q26发布，是腾讯在Agent-to-Agent领域的关键落子。混元Hy3正式版输入价格降至RMB1/百万token（预览版RMB2），任务成功率从72\%提升至90\%，token消耗节省近50\%。企业级AI助手WorkBuddy桌面月访问量达885万次，领先字节Trae和阿里Qoder。
{\small\color{refgray}数据：JEF预计腾讯2Q26非IFRS净利润增速仅4\%，远低于市场共识9.7\%\par}
{\small\color{refgray}数据：微信AI‘小微’6月内测，正式版预计4Q26发布\par}
{\small\color{refgray}数据：混元Hy3正式版输入价格降至RMB1/百万token（预览版RMB2）\par}
{\small\color{refgray}数据：混元Hy3任务成功率从72\%提升至90\%，token消耗节省近50\%\par}
{\small\color{refgray}边际变化：JEF指出腾讯AI产品从2026年才开始显著贡献收入，而资本开支从2025年已大幅上升，时间错配导致短期利润增速低于收入增速。\par}
\subsubsection*{关键数据}
\begin{itemize}
  \item BofA自建存储景气指标5月录得183，远高于历史周期高点120-130
  \item DRAM现货价格同比涨幅达700\%
  \item 南亚科技2Q26 DRAM ASP环比涨超60\%，营业利润率从61\%跳升至74\%
  \item GS推算三星电子2Q26常规DRAM OPM达80\%
  \item 晶合集成40nm和28nm扩产总预算355亿元，预计2028年二季度满产
  \item JEF预计腾讯2Q26非IFRS净利润增速仅4\%，远低于市场共识9.7\%
  \item 混元Hy3正式版任务成功率从72\%提升至90\%，token消耗节省近50\%
\end{itemize}
\begin{figure}[htbp]
\centering
\includegraphics[width=0.88\linewidth]{figures/fig_001.jpg}
\caption{Exhibit 3 - Exhibit 3: New record-high in Jun (NT\textbackslash{}\$29.4bn; +6\% MoM); more than 5x higher than 2025-avg (NT\textbackslash{}\$5.5bn) Nanya Tech – Monthly sales (Jun 2026) BofA GLOBAL RESEARCH Nanya Tech DRAM ASP trend – quarterly; 8Gb equiv. BofA GLOBAL RES}
\end{figure}
\begin{figure}[htbp]
\centering
\includegraphics[width=0.88\linewidth]{figures/fig_018.jpg}
\caption{Exhibit 3 - Exhibit 3: Peers' correlation of forward year trading P/E and net income growth Exhibit 4: Nexchip discounted P/E}
\end{figure}
\begin{figure}[htbp]
\centering
\includegraphics[width=0.88\linewidth]{figures/fig_002.jpg}
\caption{Exhibit 3 - Exhibit 7: Also at all-time high in Jun (NT\textbackslash{}\$24.9bn; +9\% MoM); more than 4x higher than 2025-avg (NT\textbackslash{}\$6.1bn) Phison Electronics – Monthly sales (Jun 2026) Exhibit 4: Robust YoY growth continued in Jun (+621\%); already 11 consecut}
\end{figure}
\begin{figure}[htbp]
\centering
\includegraphics[width=0.88\linewidth]{figures/fig_003.jpg}
\caption{Exhibit 5 - Exhibit 7: Also at all-time high in Jun (NT\textbackslash{}\$24.9bn; +9\% MoM); more than 4x higher than 2025-avg (NT\textbackslash{}\$6.1bn) Phison Electronics – Monthly sales (Jun 2026) Exhibit 4: Robust YoY growth continued in Jun (+621\%); already 11 consecut}
\end{figure}
{\small\color{refgray}本节综合 4 篇研究；涉及 BofA、GS、JEF。}
\newpage
\subsection{汽车与出行：电动化分化、经销商阵痛与机器人新叙事}
\textbf{全球汽车产业正经历结构性分化：欧洲电动化加速由平价车型驱动，中国渗透率攀升但总量承压，美国EV渗透率逆势下滑；豪华车经销商尚未触底，售后业务恶化；日本车企重新将人形机器人视为增长引擎，但自动化产业链利润修复仍需时间。}
\subsubsection*{主流共识}
\begin{itemize}
  \item 欧洲EV+PHEV渗透率在平价车型推动下显著提升，EU-5六月同比增8.6个百分点至34.5\%。
  \item 中国新能源车渗透率突破63\%，但乘用车总销量同比降23\%，高渗透率掩盖了总量调整。
  \item 豪华车经销商新车亏损收窄仅为数学效应，售后业务因门店关闭而恶化，行业整合尚未完成。
\end{itemize}
\subsubsection*{分歧与条件差异}
\begin{itemize}
  \item 欧洲电动化动力：JPM认为由平价大众市场车型驱动，而非补贴；Bernstein关注中国车企份额上升至约20\%的竞争压力。
  \item 中国新能源车前景：JPM维持2026年乘用车需求同比降15\%的判断，认为农村补贴仅为存量稳定器；市场部分观点期待政策刺激反转。
  \item 日本车企机器人战略：Bernstein认为三菱量产计划（2027年月产1000台）是商业化信号；MS指出安川电机订单回暖但利润率仅1.6\%，盈利转化存疑。
\end{itemize}
\subsubsection*{机构观点}
\paragraph{Bernstein} 日本车企正重新将人形机器人视为中长期增长引擎，三菱与Highlanders合作计划2027年实现月产1000台，利用闲置工厂降低试错成本。全球人形机器人市场预计2050年达7290亿美元，车企在电机、传感器等核心技术上有天然优势。本田、丰田虽有长期研发积累，但三菱是首家将量产提上明确日程的日本车企。
{\small\color{refgray}数据：三菱计划2027年月产1000台人形机器人\par}
{\small\color{refgray}数据：全球人形机器人市场2050年预计7290亿美元\par}
{\small\color{refgray}数据：2025-2050年出货量CAGR约37\%\par}
{\small\color{refgray}数据：本田1986年开始双足机器人研究\par}
{\small\color{refgray}边际变化：日本车企从研发展示转向明确量产规划，三菱首次设定产能目标并利用闲置产能。\par}
\paragraph{Bernstein} 大众“未来计划”缺乏执行细节，市场反应冷淡。计划将车型精简50\%、复杂度降75\%、年产能从1000万降至900万，但未说明关停工厂、时间表及成本节约目标。同日媒体曝出可能关闭茨维考、埃姆登等工厂涉及4万岗位，新闻稿却只字未提。劳资博弈与2024年谈判高度相似，重组能否推进存疑。
{\small\color{refgray}数据：车型精简最多50\%\par}
{\small\color{refgray}数据：复杂度减少最多75\%\par}
{\small\color{refgray}数据：年产能从1000万降至900万辆\par}
{\small\color{refgray}数据：媒体曝关闭工厂涉及约4万个岗位\par}
{\small\color{refgray}边际变化：市场从期待重组细节转向质疑执行力，劳资博弈历史重演风险上升。\par}
\paragraph{JPM} 全球电动化区域分化加剧：欧洲EU-5六月EV+PHEV渗透率同比升8.6个百分点至34.5\%，由平价车型驱动；中国NEV渗透率突破63\%但乘用车总量同比降23\%，农村补贴仅为存量稳定器；美国EV渗透率逆势降至6.6\%。储能系统正取代电动车成为电池产业链最确定增量，上调2030年美国储能需求预测52\%至303GWh，数据中心占45\%。
{\small\color{refgray}数据：EU-5六月EV+PHEV渗透率34.5\%，同比+8.6ppt\par}
{\small\color{refgray}数据：中国六月NEV零售100万辆，渗透率63\%\par}
{\small\color{refgray}数据：中国乘用车总销量同比-23\%\par}
{\small\color{refgray}数据：美国六月EV渗透率6.6\%，同比-3ppt\par}
{\small\color{refgray}边际变化：储能超越电动车成为电池需求主驱动，美国FERC改革将数据中心审批缩短至90天。\par}
\paragraph{MS} 中国豪华车经销商尚未触底，行业整合仍在进行。2026上半年BBA合资公司零售量同比降16\%-28\%，保时捷降32\%。新车亏损收窄仅为销量下降的数学效应，售后业务因门店关闭恶化：中升售后收入增速从2023年16\%降至2025年4\%，永达4S店从179家减至137家。社保趋严和商誉减值风险进一步压制估值修复。
{\small\color{refgray}数据：BBA中国零售量同比降16\%-28\%\par}
{\small\color{refgray}数据：保时捷中国销量降32\%\par}
{\small\color{refgray}数据：中升售后收入增速从16\%降至4\%\par}
{\small\color{refgray}数据：永达4S店从179家减至137家\par}
{\small\color{refgray}边际变化：售后业务从利润稳定器转为拖累，门店关闭导致客户流失至独立维修店。\par}
\paragraph{MS} 安川电机订单回暖但利润端承压，1季度营业利润同比降19\%至85亿日元，远低于预期。机器人订单环比增37\%超预期，但利润率降至1.6\%，受ERP系统迁移和重组成本影响。运动控制销售额未能同比增长。管理层维持全年指引，但利润修复需待2季度生产正常化。
{\small\color{refgray}数据：1季度营业利润85亿日元，同比-19\%\par}
{\small\color{refgray}数据：机器人订单环比+37\%\par}
{\small\color{refgray}数据：机器人利润率1.6\%\par}
{\small\color{refgray}数据：AC伺服订单同比+65\%\par}
{\small\color{refgray}边际变化：订单复苏尚未转化为盈利改善，ERP迁移导致生产中断和成本上升。\par}
\subsubsection*{关键数据}
\begin{itemize}
  \item EU-5六月EV+PHEV渗透率34.5\%，同比+8.6ppt
  \item 中国六月NEV渗透率63\%，但乘用车总销量同比-23\%
  \item 美国六月EV渗透率6.6\%，同比-3ppt
  \item 三菱计划2027年月产1000台人形机器人
  \item 全球人形机器人市场2050年预计7290亿美元
  \item BBA中国零售量同比降16\%-28\%
  \item 安川电机1季度营业利润85亿日元，同比-19\%
  \item 2030年美国储能需求预测上调52\%至303GWh
\end{itemize}
\begin{figure}[htbp]
\centering
\includegraphics[width=0.88\linewidth]{figures/fig_025.jpg}
\caption{Figure 1 - Figure 1: Global EV battery installations by company vs KOR battery cell market share Figure 2: Global EV battery installation market share trend by company \% Figure 3: EU10 + US + China EV/PHEV sales and penetration trend Uni}
\end{figure}
\begin{figure}[htbp]
\centering
\includegraphics[width=0.88\linewidth]{figures/fig_031.jpg}
\caption{Exhibit 1 - Exhibit 1: BBA China retail volume continues to decline YoY Exhibit 2: BMW China market share (\%) Exhibit 3: Mercedes China market share (\%)}
\end{figure}
\begin{figure}[htbp]
\centering
\includegraphics[width=0.88\linewidth]{figures/fig_006.jpg}
\caption{EXHIBIT 1 - EXHIBIT 1: Highlanders' humanoid robot 'N' EXHIBIT 2: Honda's history of robotics development EXHIBIT 3: Toyota's 'Mobility Robot'}
\end{figure}
\begin{figure}[htbp]
\centering
\includegraphics[width=0.88\linewidth]{figures/fig_026.jpg}
\caption{Figure 1 - Figure 4: EU10 + US + China EV sales and penetration trend}
\end{figure}
{\small\color{refgray}本节综合 5 篇研究；涉及 Bernstein、JPM、MS。}
\newpage
\subsection{消费与餐饮：茶饮韧性修复，正餐复苏仍待确认}
\textbf{6月中国餐饮数据延续分化：现制茶饮在高基数下展现超预期韧性，头部品牌通过产品创新和品类扩张实现同店销售环比改善；而火锅与正餐受天气及宏观复苏偏弱影响，翻台率与同店增速承压，暑期消费前景需重新校准。}
\subsubsection*{主流共识}
\begin{itemize}
  \item 现制茶饮头部品牌同店销售环比改善，韧性优于市场预期。
  \item 正餐板块（海底捞、太二）6月同店增速放缓，天气与宏观复苏偏弱是主因。
  \item 下半年基数降低，但暑期出行数据（酒店RevPAR、国内航班）暗示客流可能仍不温不火。
\end{itemize}
\subsubsection*{分歧与条件差异}
\begin{itemize}
  \item 茶饮改善是“有条件的修复”还是可持续趋势？新品创新能否持续支撑客流。
  \item 正餐品牌门店扩张（太二6.0新店、海底捞加盟店）能否对冲同店下滑压力。
  \item 促销力度是否将加大？当前克制促销与市场担忧的价格战之间存在分歧。
\end{itemize}
\subsubsection*{机构观点}
\paragraph{GS} 6月现制茶饮板块同店销售在高基数下展现韧性，优于市场担忧。茶百道同店销售环比改善至仅小幅下滑，受益于咖啡新品及“冰椰荔枝”成功；古茗同店改善至微降，部分因促销日历调整；霸王茶姬月均GMV超35万元，新品贡献约15\%GMV。促销力度整体克制，品牌通过产品创新而非单纯价格战维持客流。正餐方面，海底捞翻台率约3.8次，同比小幅下滑，恢复至2019年约80\%；太二同店增速放缓至中个位数下滑，最后一周受天气影响。太二6.0新店首周销售比老店高30\%-40\%，海底捞上半年净关8家直营店、开设14家加盟店。
{\small\color{refgray}数据：茶百道6月同店销售环比改善至仅小幅下滑\par}
{\small\color{refgray}数据：霸王茶姬月均GMV超35万元，新品贡献约15\%\par}
{\small\color{refgray}数据：海底捞6月翻台率约3.8次，恢复至2019年约80\%\par}
{\small\color{refgray}数据：太二6月同店增速中个位数下滑，前三周仍双位数增长\par}
{\small\color{refgray}边际变化：茶饮板块从市场担忧的持续下滑转为环比改善，韧性超预期；正餐板块同店增速从5月高个位数放缓至中个位数，边际走弱。\par}
\paragraph{MS} 海底捞6月翻台率同比小幅下滑，低于预期。端午假期未能扭转趋势，不利天气是部分原因，但宏观复苏偏弱是更根本的驱动。下半年基数虽降低，但酒店RevPAR和国内航班等出行数据暗示暑期客流可能依然不温不火，市场对消费复苏的期待需重新校准。当前估值按2026年预期盈利约10.5倍市盈率，低于历史均值。
{\small\color{refgray}数据：海底捞6月翻台率同比小幅下滑，低于MS预期\par}
{\small\color{refgray}数据：端午假期未能改善翻台率同比趋势\par}
{\small\color{refgray}数据：酒店RevPAR和国内航班数据暗示暑期客流偏弱\par}
{\small\color{refgray}数据：当前估值约10.5倍2026年预期市盈率\par}
{\small\color{refgray}边际变化：从5月翻台率同比小幅下降延续至6月，且端午假期未能带来改善，暑期前景从温和复苏预期转为谨慎。\par}
\subsubsection*{关键数据}
\begin{itemize}
  \item 茶百道6月同店销售环比改善至仅小幅下滑
  \item 霸王茶姬月均GMV超35万元，新品贡献约15\%
  \item 海底捞6月翻台率约3.8次，恢复至2019年约80\%
  \item 太二6月同店增速中个位数下滑，前三周仍双位数增长
  \item 太二6.0新店首周销售比老店高30\%-40\%
  \item 海底捞上半年净关8家直营店，开设14家加盟店
  \item 海底捞当前估值约10.5倍2026年预期市盈率
  \item 酒店RevPAR和国内航班数据暗示暑期客流偏弱
\end{itemize}
\begin{figure}[htbp]
\centering
\includegraphics[width=0.88\linewidth]{figures/fig_020.jpg}
\caption{Exhibit 3 - Exhibit 3: Store opening monthly tracker Haidilao brands store opening based on company data; Jiumaojiu store opening based on brand official account; others based on Zhaimen. Exhibit 4: Jiumaojiu operating results snapshot \#}
\end{figure}
\begin{figure}[htbp]
\centering
\includegraphics[width=0.88\linewidth]{figures/fig_030.jpg}
\caption{Exhibit 1 - \textbackslash{}- Unfavorable weather likely partly explains soft traffic, but weak macro recovery remains the key driver. Although the YoY base eases in 2H, travel trends (hotel RevPAR, domestic flights) suggest summer traffic may stay lukewarm. See also: Hong Kong/China Le}
\end{figure}
\begin{figure}[htbp]
\centering
\includegraphics[width=0.88\linewidth]{figures/fig_021.jpg}
\caption{Exhibit 4 - Exhibit 4: Jiumaojiu operating results snapshot \#\# High frequency indicators Exhibit 5: DAU/user session of major delivery/food service brand APP}
\end{figure}
\begin{figure}[htbp]
\centering
\includegraphics[width=0.88\linewidth]{figures/fig_022.jpg}
\caption{Exhibit 4 - Exhibit 4: Jiumaojiu operating results snapshot \#\# High frequency indicators Exhibit 5: DAU/user session of major delivery/food service brand APP}
\end{figure}
{\small\color{refgray}本节综合 2 篇研究；涉及 GS、MS。}
\newpage
\subsection{区域市场与策略：结构性分化下的机会与风险}
\textbf{当前全球市场呈现显著的结构性分化：中国电子材料周期强劲但粗纱需求疲软，印度市场在外资抛售后迎来反弹窗口，而油轮运输则受地缘事件扰动但基本面稳健。投资者需关注行业间的错位机会。}
\subsubsection*{主流共识}
\begin{itemize}
  \item 中国电子纱/电子布供需紧张，库存低位，价格维持高位，下游覆铜板需求旺盛。
  \item 印度市场外资持仓已降至极低水平，资金回流空间充足，大盘股相对中盘股折价明显。
  \item 油轮运输有效运力受地缘事件稀释，但核心航线运价同比大幅上涨，行业景气度持续。
\end{itemize}
\subsubsection*{分歧与条件差异}
\begin{itemize}
  \item 中国粗纱与电子纱走势分化：粗纱受建筑淡季影响价格承压，而电子纱受益于AI和高端消费电子需求。
  \item 印度盈利下修周期是否结束：部分行业（能源、旅游）可能被过度修正，而金属、医疗下修可能不足。
  \item 油轮运价上涨的可持续性：二季度环比改善有限，三季度有效利用率恢复后运价能否维持是关键。
\end{itemize}
\subsubsection*{机构观点}
\paragraph{Citi} 中国玻璃纤维市场结构性分化明显：粗纱价格在季节性淡季中承压，2400tex直接纱均价3778元/吨，周环比微降0.33\%；而电子纱G75主流价16800-17500元/吨，周环比上涨3.6\%，电子布7628维持8.5-8.7元/米高位。新增产能爬坡但有效供给有限，电子纱/布库存处历史低位，下游覆铜板客户积极锁货，部分小客户需加价获取配额。电子材料周期尚未见顶，近零库存和持续供需紧张构成核心支撑。
{\small\color{refgray}数据：2400tex直接纱均价3778元/吨，周环比-0.33\%\par}
{\small\color{refgray}数据：G75电子纱主流价16800-17500元/吨，周环比+3.6\%\par}
{\small\color{refgray}数据：7628电子布维持8.5-8.7元/米高位\par}
{\small\color{refgray}数据：在产窑炉122座，总产能903.8万吨/年\par}
{\small\color{refgray}边际变化：粗纱需求季节性减弱，电子纱供需紧张加剧，下游抢货行为增多。\par}
\paragraph{GS} 印度市场回调后底部已形成，NIFTY指数至2027年6月目标26500点，较当前约有10\%上行空间。外资上半年创纪录卖出约300亿美元后，6月中旬已转为净买入20亿美元，主要集中于金融板块。盈利下修周期接近尾声，MSCI印度2026年盈利增长预期从16\%下修至12\%，接近GS全年10\%预测，但能源精炼、旅游航空等板块可能被过度修正。大盘股估值回落至15年均值附近，相对中盘股折价达30\%，提供更好盈利可见性。
{\small\color{refgray}数据：NIFTY指数上半年回调9\%\par}
{\small\color{refgray}数据：外资上半年净卖出约300亿美元，创历史纪录\par}
{\small\color{refgray}数据：MSCI印度2026年盈利增长预期从16\%下修至12\%\par}
{\small\color{refgray}数据：大盘股相对中盘股折价达30\%\par}
{\small\color{refgray}边际变化：外资抛售不确定性释放完毕，资金从成长向价值轮动，大盘股相对中盘股吸引力上升。\par}
\paragraph{MS} 中远海能1H26归母净利润45亿元，同比增141\%，其中2Q26净利润约23.27亿元，环比增7\%，比预期高6\%。利润超预期主要来自油轮运价强劲，TD15（西非-中国）和TD22（美湾-中国）平均TCE同比分别上涨180\%和170\%。但环比改善有限，因部分船舶受霍尔木兹海峡事件被困，有效运力被稀释，且国际成品油和国内油运价格偏软。受影响船舶已在6月底安全驶出，三季度有效利用率有望改善。
{\small\color{refgray}数据：1H26归母净利润45亿元，同比+141\%\par}
{\small\color{refgray}数据：2Q26净利润23.27亿元，环比+7\%，比预期高6\%\par}
{\small\color{refgray}数据：TD15平均TCE同比+180\%\par}
{\small\color{refgray}数据：TD22平均TCE同比+170\%\par}
{\small\color{refgray}边际变化：二季度利润超预期但环比增幅有限，三季度有效利用率恢复后运价能否维持是关键。\par}
\subsubsection*{关键数据}
\begin{itemize}
  \item 中国G75电子纱周环比+3.6\%至16800-17500元/吨
  \item 印度NIFTY上半年回调9\%，外资创纪录卖出300亿美元
  \item 中远海能2Q26净利润23.27亿元，环比+7\%
  \item TD15和TD22平均TCE同比分别+180\%和+170\%
  \item MSCI印度2026年盈利增长预期从16\%下修至12\%
  \item 大盘股相对中盘股折价达30\%
\end{itemize}
\begin{figure}[htbp]
\centering
\includegraphics[width=0.88\linewidth]{figures/fig_011.jpg}
\caption{Figure 1 - Figure 1. Electronic yarn prices © 2026 Citi Inc. No redistribution without Citi's written permission. Figure 2. E-fabric prices © 2026 Citi Inc. No redistribution without Citi's written permission.}
\end{figure}
\begin{figure}[htbp]
\centering
\includegraphics[width=0.88\linewidth]{figures/fig_013.jpg}
\caption{Exhibit 2 - Exhibit 2: Consensus is expecting healthy 12\% yoy profit growth for MSCI India ex-commodities in 2Q CY26E Exhibit 3: Current EPS downgrade cycle in India is on track to be among the weakest in history}
\end{figure}
\begin{figure}[htbp]
\centering
\includegraphics[width=0.88\linewidth]{figures/fig_016.jpg}
\caption{Exhibit 5 - Exhibit 5: We expect consensus to cut earnings growth further by 2pp to 10\% for CY26 (from 12\%) Exhibit 6: Comparing 2026 oil shock to past episodes, recent EPS cuts may have been overdone in energy and consumer retail/services,}
\end{figure}
\begin{figure}[htbp]
\centering
\includegraphics[width=0.88\linewidth]{figures/fig_012.jpg}
\caption{Figure 1 - Figure 1. Electronic yarn prices © 2026 Citi Inc. No redistribution without Citi's written permission. Figure 2. E-fabric prices © 2026 Citi Inc. No redistribution without Citi's written permission. \#\# China Jushi (600176.SS; Rmb55.7; 1; 10 Jul 26; 15:00) \#\# V}
\end{figure}
{\small\color{refgray}本节综合 3 篇研究；涉及 Citi、GS、MS。}
\newpage
\section{辅助信号｜战略咨询}
本板块聚焦战略咨询机构对宏观产业趋势的独立分析，本期以沙特超级城市项目为案例，探讨分布式光伏如何成为新城开发的底层低碳架构，并分析零资本支出模式下的财务可行性与政策含义。
\subsection{分布式光伏：新城开发的低碳基础设施标配}
\textbf{在沙特Vision 2030框架下，超级城市项目无需等待国家电网清洁化，通过系统性利用屋顶、车棚等建筑表面安装光伏，即可覆盖高达35\%的总用电需求，实现成本对冲与碳减排。}
\begin{itemize}
  \item 典型独栋别墅屋顶光伏可满足约50\%年用电需求，中层建筑约15\%；城市开发项目组合中系统性利用所有可用表面，光伏覆盖可达35\%。
  \item 通过购电协议（PPA）或能源即服务模式，开发商可实现零资本支出，第三方负责建设运营，电价低于电网约30\%。
  \item 自建自营模式25年净现值（NPV）可达2亿-2.5亿沙特里亚尔，比PPA模式高35\%-50\%，但PPA模式零前期投入且转移运营风险。
  \item 边际变化：光伏度电成本已低于沙特居民和商业电价，开发商决策从“是否划算”转向“选择哪种模式”，回报门槛消失。
  \item 政策含义：新城规划应尽早预留光伏空间，避免后期改造成本，开发商可参考BCG提出的分阶段实施路线图。
\end{itemize}
\begin{figure}[htbp]
\centering
\includegraphics[width=0.88\linewidth]{figures/fig_107.jpg}
\caption{Exhibit 1 - Power purchase agreements (PPAs) and energy-as-a-service contracts have effectively removed the barrier of requiring upfront investment. These options enable solar systems to be installed and operated with no upfront cost, allowing developers to buy electricit}
\end{figure}
\newpage
\section{辅助信号｜智库/国际机构}
本板块整合了亚洲开发银行、布鲁盖尔研究所、IMF、经合组织和世界贸易组织的最新研究，聚焦三大主题：中国债市基础设施兼容性、AI收益分配的结构性不对称、以及全球供应链与贸易规则的深层变化。
\subsection{中国债市开放：基础设施兼容性决定资金长期停留}
\textbf{亚洲开发银行指出，Bond Connect已从市场准入选项演变为覆盖交易、托管、外汇、衍生品和回购的完整跨境生态，其核心价值在于让国际投资者用熟悉的平台和习惯无缝接入，而非简单的额度放开。}
\begin{itemize}
  \item Bond Connect允许境外投资者参与CIBM买断式回购业务（2025年），并通过Swap Connect提供利率互换工具，覆盖交易、结算、对冲和融资全环节。
  \item 与QFII/RQFII和CIBM Direct相比，Bond Connect的核心优势是操作兼容性——投资者可继续使用Bloomberg、MarketAxess、Tradeweb等全球交易平台。
  \item 报告强调，判断市场是否真正开放，不能只看政策文件中的“允许”，更要看国际投资者能否以可接受的成本完成交易、托管和对冲。
\end{itemize}
\begin{figure}[htbp]
\centering
\includegraphics[width=0.88\linewidth]{figures/fig_036.jpg}
\caption{Figure 3 - Figure 3: Breakdown of Northbound Bond Connect Investors by Jurisdiction Figure 4: Breakdown of Northbound Bond Connect Investors by Investor Type Figure 5: Assets Under Custody by International Institutional Investors in the}
\end{figure}
\subsection{AI收益分配的结构性不对称：2.7万亿美元时间价值集中于高收入国家}
\textbf{IMF最新工作论文估算，当前AI使用带来的时间节省价值约为每年2.7万亿美元（占样本国家GDP的3.4\%），但低收入国家几乎未受益，AI收益分配正经历从“窄到宽”的转变，但速度高度依赖收入水平、监管准备度和语言环境。}
\begin{itemize}
  \item AI节省的时间价值按各国工资水平估算（劳动成本等价LCE），低收入国家因工资低且AI对话集中在极窄职业范围，贡献份额极小。
  \item 美国AI使用广度从软件开发扩散到教育、销售、办公室工作，驱动LCE增长；而坦桑尼亚的AI使用几乎全部集中在专业职业领域，农业、服务业等占劳动力大多数的群体几乎不参与。
  \item 非英语国家的AI使用数据若出现职业扩散，可能意味着多语言模型能力的实质提升，这是发展中国家能否从AI中获益的关键检验。
\end{itemize}
\begin{figure}[htbp]
\centering
\includegraphics[width=0.88\linewidth]{figures/fig_045.jpg}
\caption{Figure 2 - Figure 2: Aggregate AI gains: scale, trend, and incidence by income. Panel (a) (R1 to R5, Jan 2025 to Feb 2026): the blue line (left axis) is the wage index, the average US occupation-level wage weighted by global AI conversation}
\end{figure}
\subsection{全球供应链骨架未动摇，但规则与竞争格局正在重塑}
\textbf{经合组织与WTO的最新研究显示，全球价值链贸易量（剔除通胀）占GDP比重仍维持在17\%历史高位，并未系统性下降；同时，WTO将化石燃料补贴改革推进至实操阶段，欧盟产业加速法案的产地规则可能扭曲资源配置。}
\begin{itemize}
  \item 经合组织首次提供“上年价格”投入产出表，剥离价格效应后，2024年全球价值链贸易量占比与2022年高点持平，企业并未大规模放弃跨境采购。
  \item WTO部长级会议将化石燃料补贴改革从联合声明推进至产出文件阶段，建立透明化提问清单和扰动补贴设计指南，要求补贴透明、精准、临时。
  \item 布鲁盖尔研究所指出，欧盟《产业加速法案》设2035年制造业GDP占比20\%的硬目标，缺乏宏观环境逻辑，可能引导资金流向政治优先级而非市场效率最高的领域，且“欧洲制造”产地规则将推高成本、延迟脱碳并面临WTO挑战。
\end{itemize}
\begin{figure}[htbp]
\centering
\includegraphics[width=0.88\linewidth]{figures/fig_102.jpg}
\caption{Figure 2 - Before the global financial crisis, both nominal and real measures show a strong increase in trade via GVCs, confirming their rapid expansion during the period of “hyperglobalisation” (Subramanian and Kessler, 2014[4]). After 2011, however, the two series dive}
\end{figure}
\begin{figure}[htbp]
\centering
\includegraphics[width=0.88\linewidth]{figures/fig_043.jpg}
\caption{Figure 1 - Figure 1: European battery cell manufacturing capacity by company, region of HQ and operational status (as of Q1 2026) FDI screening has historically been a national competence. The EU introduced only in 2019 coordination of scr}
\end{figure}
\subsection{宏观政策与金融稳定}
\textbf{用于校准利率、通胀、财政约束和金融稳定叙事。}
\begin{itemize}
  \item 国际研究争端解决正在经历一场静默的转向。仲裁成本飙升、周期漫长，调解作为一种低成本、高效率的替代方案，理论上备受推崇。但亚洲开发银行研究所最新政策简报指出，调解在读者与国家间的争端中远未普及，核心障碍并非技术或法律空白，而是政府自身的运作逻辑。
  \item 这份来自国际货币产品组织（IMF）的研报，针对纳米比亚的社会支出提出了一个反直觉的判断：这个国家在教育、医疗和社会救助上的公共支出水平，放在同等发展水平的国家中并不低，甚至偏高。但真正的问题在于，这些钱花下去之后，并没有带来同等的产出。
  \item 一份刚刚发布的IMF工作论文，用2002至2021年间的企业级数据，对中东和中亚（ME\&CA）地区的市场竞争格局做了迄今最系统的量化分析。结论直接、有政策分量：这个地区的生产率增长瓶颈，不在资源禀赋，而在竞争不足。而关税，是其中最可干预的变量。
  \item 一份来自IMF技术援助团队的最新评估报告，对多米尼加共和国的财政透明度做了全面诊断。结论是：这个国家的财政报告和预算流程已有相当基础，但真正的不确定性藏在报表之外——公共信托产品、长期购电协议、养老金缺口和国企财务状况，这些关键领域要么数据缺失，要么分析流于形式。报告的核心判断是...
  \item 一份2025年5月完成、2026年1月发布的IMF技术援助报告，揭示了一个多数市场观察者容易忽略的事实：在非洲南部小国史瓦帝尼，非银行机构的资产规模已占据整个资金体系的66.2\%，但此前长期游离在官方货币与资金统计框架之外。报告的核心判断是，IMF正在帮助该国央行将养老金、保险、...
  \item 伊拉克宏观环境长期被一个问题困扰：高度依赖石油收入，油价波动直接冲击财政、货币和国际收支。中央银行需要一套能预测这些冲击、支持政策决策的分析框架。IMF的技术援助报告给出的结论很直接——工具已经有了，真正的挑战不在模型本身，而在团队能否持续用它、数据能否支撑它、央行能否与财政部真...
  \item 在固定汇率制度下，货币政策独立性天然受到制约。但这份IMF最新工作论文揭示了一个更具体的结论：纳米比亚的银行存贷款利率变动，主导力量来自南非储备银行的政策利率，而非本国央行的回购利率。这不仅是学术验证，更对任何实行货币局制度或紧密盯住汇率的地区，提供了一个关于政策传导边界的现实参...
  \item 一份来自国际货币产品组织（IMF）的最新工作论文，系统拆解了美国回购市场利差的驱动因素。结论并不复杂，但含义值得细品：准备金依然是回购市场最核心的稳定器，但它的效果并非在所有市场环境下都相同。真正决定回购利率走向的，是准备金、交易商资产负债表使用率、对冲产品杠杆与国债发行这四个变...
  \item 2025年3月，联合国统计委员会与国际货币产品组织同步发布了国民账户体系和国际收支手册的最新版本。对大多数国家的统计机构而言，这不仅仅是一次技术升级，而是一场需要跨部门协调、用户参与和资源重新配置的系统工程。IMF统计部门最新发布的《SNA/BPM实施规划手册》给出了一个清晰的信...
  \item 一份来自国际货币产品组织（IMF）的最新工作论文，为理解稳定币在新兴市场的角色提供了一个意想不到的分析框架。报告的核心判断是：稳定币既是改善外汇准入和价格发现的技术工具，也可能成为协调资本外逃的“公共信号放大器”——其福利效应取决于汇率失衡的程度。
  \item 理解资金体系之外的宏观环境活动，对政策制定者而言从来不是可选项。但加密货币挖矿——这个不依赖传统资金中介、直接创造新资产的入口——其真实规模在全球范围内几乎是一团迷雾。IMF最新工作论文提出了一种全新的测量思路：通过追踪全球主要产地的矿机出口数据，来反推各国的挖矿活动分布。这个方...
  \item 全球固定收益市场在2024年达到145.1万亿美元规模，而同期股权市场为126.7万亿美元。主权和企业的借贷还在攀升。在这片汪洋中，绿色、社会、可持续和可持续发展挂钩债券（GSSS债券）被寄予厚望——它们能够将规模、可持续影响和长期久期结合在一起。经合组织与卢森堡标的交易所联合发...
\end{itemize}
\begin{figure}[htbp]
\centering
\includegraphics[width=0.88\linewidth]{figures/fig_050.jpg}
\caption{Figure 1 - 4. The paper is organized as follows. Section B provides the conceptual framework and situates social spending within the broader fiscal context. Section C assesses the education sector. Section D examines the health sector. Section E analyzes social protectio}
\end{figure}
\begin{figure}[htbp]
\centering
\includegraphics[width=0.88\linewidth]{figures/fig_055.jpg}
\caption{Figure 1 - \#\# 4.4 Stylized Facts \#\# 4.4.1 Market Power in ME\&CA Figure 1. Sales-weighted Markups in ME\&CA and Other EMDEs Sources: Orbis database by Bureau van Dijk, and the authors' estimates. Turning to country-level dynamics, the data reveal meaningful heterogeneity a}
\end{figure}
\subsection{AI、技术与生产率}
\textbf{用于判断 AI 投资周期、生产率扩散和产业约束是否正在改变资产定价。}
\begin{itemize}
  \item 一份来自亚洲开发银行的工作论文，把目光投向了全球最被低估的住房挑战之一：兴都库什-喜马拉雅地区。这片“亚洲水塔”支撑着超过十亿人的生计，但其快速城市化进程正与极高的地震不确定性正面碰撞。报告的核心判断是，这三个国家（不丹、印度喜马拉雅地区、尼泊尔）的住房市场存在一个系统性的“脱节...
  \item 一份新的IMF技术援助报告揭示了一个关键进展：伊拉克中央银行（CBI）已成功开发并掌握了宏观宏观环境基础工具（MFT），但该工具要真正融入政策决策过程，还需跨越数据协调与内部流程的最后障碍。
  \item 一份由经合组织（OECD）为马耳他撰写的技能战略报告，看起来像是一份标准的小国政策建议。但仔细读完30条建议和背后的诊断，你会发现一个反直觉的判断：马耳他技能体系的核心瓶颈，不是资金不足，不是学校太少，不是培训机会不够，而是“谁在负责”这件事从未被真正回答。
  \item 过去两年，AI模型的性能在快速提升，价格在急剧下降。经合组织最新发布的报告确认了这一趋势：其构建的AI模型质量调整价格指数，从2024年1月到2026年4月下降了近80\%。但这份报告的核心判断并不乐观——它认为，这些短期动态可能掩盖了更深层的结构性不确定性。真正的问题不是今天AI...
\end{itemize}
\begin{figure}[htbp]
\centering
\includegraphics[width=0.88\linewidth]{figures/fig_039.jpg}
\caption{Figure 1 - Over the last 3 decades (1992–2022), the Global Emergency Management Database, Emergency Events Database (EM-DAT), recorded 264 disasters in the study area, leading to fatalities (71,691 deaths), displacement (1.2 billion people affected, including 8.4 million}
\end{figure}
\begin{figure}[htbp]
\centering
\includegraphics[width=0.88\linewidth]{figures/fig_060.jpg}
\caption{Figure 1 - Once the forecast discussed and finalized, prepare a draft forecast presentation. \textbackslash{}- Finally, prepare few alternative scenarios which usually focus on different oil price trajectories or fiscal adjustments or any other major macro shock facing Iraqi economy. T}
\end{figure}
\subsection{地缘政治与安全}
\textbf{用于补充能源、航运、防务和供应链风险的尾部情景。}
\begin{itemize}
  \item 全球水安全研究的回报率已被反复证明——联合国估算，每投入1美元，平均能产生约4美元的收益。但现实是，全球水相关基础设施长期面临资金不足、资产过早老化、服务交付效率低下等问题。经合组织（OECD）最新发布的《水安全融资》报告给出了一个核心判断：问题不在于缺钱，而在于钱怎么投、怎么回...
  \item 一份来自经合组织的最新跟进报告揭示了土耳其反海外贿赂体系中的一个尖锐矛盾：政府持续做出改革承诺，但实际执行进展却几近停滞。在2024年第四轮评估提出的71项建议中，土耳其仅完全落实了10项，仍有49项未采取任何实质性行动。这是该工作组自2007年以来反复敦促却始终未见突破的领域。
\end{itemize}
\begin{figure}[htbp]
\centering
\includegraphics[width=0.88\linewidth]{figures/fig_092.jpg}
\caption{Figure 2 - Unlike GSS bonds, which allocate proceeds to specific projects, “sustainability-linked” bonds (SLBs) are general-purpose bonds, whose financial terms are tied to the achievement of environmental and social performance targets (OECD, 2024[9]). In other words, t}
\end{figure}
\section{结语}
后续需验证的关键节点包括：存储芯片三季度出货量及价格走势是否确认结构性分化；中国新能源车总量能否在政策刺激下企稳；日本车企机器人订单向利润的转化进度；以及油轮运价在有效利用率恢复后的实际表现。
\newpage
\section{覆盖概览}
投行/券商 14 篇；战略咨询 1 篇；智库/国际机构 23 篇。\par
投行覆盖：GS 4篇 / MS 4篇 / Bernstein 2篇 / BofA 1篇 / Citi 1篇 / JEF 1篇 / JPM 1篇\par
\section{Disclaimer}
{\scriptsize\color{refgray}
This community is a paid learning and information-sharing community created for general educational purposes only. The membership fee is charged solely for access to community discussions, educational content, organization, and operational costs. It is not an advisory fee, management fee, performance fee, brokerage fee, or compensation for personalized financial, investment, legal, tax, or professional advice.\par
I participate and share information solely as an individual and for educational discussion among community members. I am not acting as a financial adviser, investment adviser, broker, dealer, portfolio manager, legal adviser, tax adviser, accountant, fiduciary, or any other licensed professional adviser. Nothing shared in this community constitutes, and should not be interpreted as, financial advice, investment advice, legal advice, tax advice, a recommendation, solicitation, offer, endorsement, instruction, or guarantee to buy, sell, hold, short, trade, or invest in any security, cryptocurrency, fund, derivative, strategy, or other financial product.\par
All content shared in this community, including messages, discussions, links, files, charts, examples, market commentary, watchlists, AI-generated or AI-polished summaries, and third-party materials, is general, impersonal, and educational in nature. It does not take into account any individual's financial situation, investment objectives, risk tolerance, investment horizon, liquidity needs, tax status, legal restrictions, or suitability requirements. No content should be relied upon as suitable for any specific person, account, portfolio, or financial decision.\par
Information shared in this community may be incomplete, inaccurate, outdated, speculative, AI-generated, AI-edited, or based on third-party internet sources that have not been independently verified. I make no representation or warranty as to the accuracy, completeness, timeliness, reliability, or suitability of any information shared.\par
Financial markets involve substantial risk, including the possible loss of principal. Past performance, historical data, hypothetical examples, simulated results, backtests, personal opinions, or educational examples do not guarantee future results. Each member is solely responsible for conducting their own research, exercising independent judgment, and making their own decisions. Before making any financial, investment, legal, or tax decision, members should consult qualified and licensed professionals.\par
Participation in this community, payment of any membership fee, or communication with me or other members does not create any adviser-client, fiduciary, professional, contractual, agency, partnership, employment, or other advisory relationship. By joining or participating in this community, you acknowledge and agree that you are solely responsible for your own decisions, actions, transactions, profits, losses, risks, and consequences, and that I shall not be liable for any direct, indirect, incidental, consequential, financial, legal, tax, or other loss or damage arising from or related to any content, discussion, information, or material shared in this community.\par
}
\newpage
\begin{center}
{\Large 更多详情报告kcdesk.com}\par\vspace{0.8cm}
\includegraphics[width=0.58\linewidth]{\detokenize{../../prompts/zsxq_img.jpg}}
\end{center}
\end{document}
